\documentclass[10pt,a4paper]{amsart}
\usepackage[margin=19mm]{geometry}
\usepackage{amsmath,amssymb,mathtools}
\usepackage[scr=rsfs,cal=euler]{mathalfa}
\usepackage{xfrac}
\usepackage[unicode,hidelinks,bookmarks=false]{hyperref}
\newcommand{\ie}{i.e.\ }
\newcommand{\cf}{cf.\ }
\newcommand{\resp}{respectively\ }
\newcommand{\Subsection}{\subsection}
\providecommand{\nobreakdash}{\nobreak\hbox{-}}
\setlength{\emergencystretch}{2em}
\multlinegap0pt
\usepackage{amssymb}
\usepackage{float}
\newtheorem{lem}{Lemma}
\title{Asymptotics and inequalities for the broken $k$-diamond partition function}
\author{Ying Zhong}
\date{Source: arXiv:2605.07502v2 (CC BY 4.0)}
\begin{document}
\maketitle

\noindent\emph{Source and licence.} Asymptotics and inequalities for the broken $k$-diamond partition function. Version arXiv:2605.07502v2 (CC BY 4.0), distributed by arXiv under the Creative Commons Attribution 4.0 International licence, http://creativecommons.org/licenses/by/4.0/, as recorded in the arXiv licence metadata for this exact version and shown on the abstract page https://arxiv.org/abs/2605.07502v2. Full licence notice: \texttt{licenses/arxiv-2605.07502-CC-BY-4.0.txt}.\par\smallskip

\noindent\textit{Editorial scope.} Counted unit: the article's lem lem:analytic (source0.tex lines 1169--1178). The article's complete original proof of it is delivered with it (source0.tex lines 1180--1266, one \texttt{proof} environment of 87 lines). No mathematics is written, repaired or re-derived: the statement and the proof are the article's own lines, in the article's own notation, and no retained line is substituted or re-flowed.  Retained uncounted as the source-local context the proof needs: lines 856--878; lines 963--969; lines 1009--1016.  Nothing was omitted from any retained span, so the package carries no omission record.  No retained line contains a citation, so the wrapper carries no bibliography and no bibliography entry.  No retained line contains a figure environment, an image-inclusion command or drawing code, so no image is delivered and the package declares no asset dependency.  Editorial formatting only, in addition: an AMS \texttt{amsart} preamble that restates the article's own declarations and macros used by the retained lines, namely the environments \texttt{lem} on separate counters, together with the standard packages those lines need.  The retained environments print in this document as Lemma 1 in order of appearance, whereas the article prints the same items with the numbering of its own full document.  The complete native source of the article as distributed in the arXiv e-print for this exact version is bundled unchanged as \texttt{sources/200086/source0.tex}.
\begin{lem}\label{lem:factor}

Assume $\varepsilon_{n-1},\varepsilon_n,\varepsilon_{n+1}<1$. Then
%\begin{equation}\label{eq:factor-core}
%\Theta(n)\le\Theta_M(n)\cdot
%\frac{(1+\varepsilon_{n-1})(1+\varepsilon_{n+1})}{(1-\varepsilon_n)^2}
%\le
%\Theta_M(n)\cdot \frac{(1+\varepsilon_n^\ast)^2}{(1-\varepsilon_n^\ast)^2},
%\end{equation}
\begin{equation}\label{eq:factor-consequence-exact}
\Theta(n)\le\Theta_M(n)\left(1+\frac{4\varepsilon_n^\ast}
{(1-\varepsilon_n^\ast)^2}\right),
\end{equation}
where $\varepsilon_n^\ast:=\max\{\varepsilon_{n-1},\varepsilon_n,\varepsilon_{n+1}\}$.
In particular, a sufficient condition for $\Theta(n)\le 1$ is
\begin{equation}\label{eq:gap-vs-error-exact}
1-\Theta_M(n)\ge\frac{4\varepsilon_n^\ast}{(1-\varepsilon_n^\ast)^2}.
\end{equation}
Moreover, if $\varepsilon_n^\ast\le \tfrac12$, it suffices that
\begin{equation}\label{eq:gap-vs-error-crude}
1-\Theta_M(n)\ge16\varepsilon_n^\ast.
\end{equation}
\end{lem}

\begin{lem}\label{lem:relerr}
Let $\delta_k:=\sqrt{\alpha_k}-\sqrt{\beta_k}>0$. For $\sqrt{\alpha_k}\,x_k(n)\ge4$, we have
\begin{equation}\label{eq:eps}
\varepsilon_n \le C_k\,x_k(n)\,e^{-\delta_k\,x_k(n)},\qquad
C_k:=\frac{432\sqrt{2}k^3 e^{\pi\sqrt{k/3}}}{\pi^{\frac{5}{2}}\alpha_k^{\frac{3}{4}}}.
\end{equation}
\end{lem}

\begin{lem}\label{lem:concavity}
For all $n$ with $x_k(n-1)\ge \dfrac{26}{\sqrt{\alpha_k}}$
%(in particular, $t(n-1):=\sqrt{\alpha_k}\,x_k(n-1)\ge 26$)
, we have
\begin{equation}\label{eq:gap}
1-\Theta_M(n) \ge \frac{\sqrt{\alpha_k}\,\pi^4}{36\,x_k^3(n+1)}.
\end{equation}
\end{lem}

\begin{lem}\label{lem:analytic}
Define
$A_k:=\frac{1728\,C_k}{\pi^4\sqrt{\alpha_k}}$.
If $x_k(n-1)\ge \dfrac{26}{\sqrt{\alpha_k}}$ and
\begin{equation}\label{eq:key-ineq}
e^{\delta_k\,x_k(n)}\ \ge\ A_k x_k^4(n)
\qquad\big(\text{equivalently }~ \delta_k x_k(n)\ge \log A_k+ 4\log x_k(n)\big),
\end{equation}
then $\Delta_k^2(n)\ge \Delta_k(n-1)\Delta_k(n+1)$.
\end{lem}

\begin{proof}
%By Lemma \ref{lem:concavity}, for $x_k(n-1)\ge 26/\sqrt{\alpha_k}$,
%\begin{equation}\label{eq:concavity-gap-36}
%1-\Theta_M(n)\ \ge\ \frac{\sqrt{\alpha_k}\,\pi^4}{36\,x_k(n+1)^3}.
%\end{equation}
By Lemma \ref{lem:relerr}, for $\sqrt{\alpha_k}\,x_k(j)\ge 4$ we have
\begin{equation*}
\varepsilon_j\ =\ \frac{|R_k(j)|}{M_k(j)}\ \le\ C_k\,x_k(j)\,e^{-\delta_k\,x_k(j)}\qquad(j=n-1,n,n+1).
\end{equation*}
Since $x_k(n)$ is increasing in $n$, we have $x_k(n-1)\le x_k(n)\le x_k(n+1)$ and therefore
\begin{equation}\label{eq:eps-star-bound}
\varepsilon_n^\ast=\max\{\varepsilon_{n-1},\varepsilon_n,\varepsilon_{n+1}\}
\le C_k\, x_k(n+1)\, e^{-\delta_k x_k(n-1)}.
\end{equation}
Using $x_k(n)=\frac{\pi}{6}\sqrt{24n-(2k+2)}$, we have
\[
x_k(n\pm1)=x_k(n)\sqrt{1\pm\frac{24}{24n-(2k+2)}}=x_k(n)\sqrt{1\pm\frac{2\pi^2}
{3x_k^2(n)}},
\]
hence
\[
x_k(n)-x_k(n-1)=x_k(n)\Big(1-\sqrt{1-\tfrac{2\pi^2}{3x_k^2(n)}}\Big)\le \frac{2\pi^2}{3x_k(n)}.
\]
Therefore,
\[
e^{\delta_k(x_k(n)-x_k(n-1))}\ \le\ \exp\!\Big(\frac{2\delta_k\pi^2}{3x_k(n)}\Big)
\ \le\ \exp\!\Big(\frac{2\sqrt{\alpha_k}\,\pi^2}{3x_k(n)}\Big)
\ \le\ \exp\!\Big(\frac{\alpha_k\pi^2}{39}\Big)=:\rho_k
\]
and
\[x_k(n+1)\le x_k(n)\sqrt{1+\frac{2\pi^2}
{3}\cdot \frac{\alpha_k}{26^2}}< \frac{3}{2}x_k(n), \]
where we use that $\delta_k\le \sqrt{\alpha_k}$, $x_k(n)\ge 26/\sqrt{\alpha_k}$ and $\alpha_k<\frac{5}{2}$. These together with \eqref{eq:eps-star-bound} give that
\begin{equation}\label{eq:eps-star}
\varepsilon_n^\ast\le C_k\, x_k(n+1)\, e^{\delta_k (x_k(n)-x_k(n-1)-x_k(n))} < \frac{3}{2}\,\rho_k\,C_k\, x_k(n)\, e^{-\delta_k x_k(n)}.
\end{equation}

Assume \eqref{eq:key-ineq}. Then by \eqref{eq:eps-star} and $x_k(n)\ge 26/\sqrt{\alpha_k}$,
\[
\varepsilon_n^\ast < \frac{3\rho_k\,C_k}{2A_k\,x_k^3(n)}
=\frac{\pi^4\rho_k\sqrt{\alpha_k}}{1152}\cdot\frac{1}{x_k^3(n)}< \frac12.
\]
so \(0\le\varepsilon_n^\ast<\tfrac12\).
%This lets us use the simplified denominator in Lemma~\ref{lem:factor}.

Then by Lemma~\ref{lem:factor}, a sufficient condition for
\(
\Theta(n)\le 1
\)
is
\begin{equation}\label{eq:sufficient-factor}
1-\Theta_M(n) \ge 16\,\varepsilon_n^\ast.
\end{equation}
%Since \(\varepsilon_n^\ast\le \tfrac12\), we have \((1-\varepsilon_n^\ast)^2\ge \tfrac14\), and hence
%\[
%\frac{4\,\varepsilon_n^\ast}{(1-\varepsilon_n^\ast)^2}\ \le\ 16\,\varepsilon_n^\ast
%\ \le\ 16\,c_\alpha\,K_k\,e^{-\delta_k x(n)}\qquad\text{by \eqref{eq:eps-star}}.
%\]
According to \eqref{eq:eps-star}, to prove \eqref{eq:sufficient-factor}, it suffices to show
\begin{equation}\label{eq:need-to-show}
1-\Theta_M(n) \ge 24\,\rho_k\,C_k\, x_k(n)\, e^{-\delta_k x_k(n)}.
\end{equation}
Using Lemma \ref{lem:concavity}, it is enough that
\[
\frac{\sqrt{\alpha_k}\,\pi^4}{36\,x_k^3(n+1)} \ge 24\,\rho_k\,C_k\, x_k(n)\, e^{-\delta_k x_k(n)},
\]
i.e.,
\[
e^{\delta_k x(n)}\ge \frac{864\,\rho_k\,C_k}{\pi^4\sqrt{\alpha_k}}\;x_k(n)x_k^3(n+1).
\]
Moreover, since
\[
\frac{x_k^3(n+1)}{x_k^3(n)}
=\Bigl(1+\frac{2\pi^2}{3x_k^2(n)}\Bigr)^{\!3/2}
\le \exp\Bigl(\frac{\pi^2}{x_k^2(n)}\Bigr)
\le \exp\Bigl(\frac{\pi^2\alpha_k}{676}\Bigr)=:\tilde{c}_k,
\]
we have
\[
\frac{864\,\rho_k\,C_k}{\pi^4\sqrt{\alpha_k}}\;x_k(n)x_k^3(n+1) \le \frac{864\,\tilde{c}_k\,\rho_k\,C_k}{\pi^4\sqrt{\alpha_k}}\;x_k^4(n).
\]
Since \(\tilde{c}_k\rho_k<2\), the right-hand side is $\le \dfrac{1728\,C_k}{\pi^4\sqrt{\alpha_k}}\,x_k^4(n)=A_k x_k^4(n)$, which is exactly the hypothesis \eqref{eq:key-ineq}. Hence \eqref{eq:need-to-show} holds, so \eqref{eq:sufficient-factor} holds, and therefore \(\Theta(n)\le 1\), i.e.,
\[
\Delta_k^2(n) \ge \Delta_k(n-1)\,\Delta_k(n+1).
\]
This concludes the proof.
\end{proof}

\end{document}
